\cotefile{45}
\scanpage{18}
% source line 1; boxed heading
\noindent\fbox{Borel \unsure{est} son propre normalisateur}

% source lines 2--5
\diagram{\begin{tikzcd}[ampersand replacement=\&,column sep=small]
\text{(i) }\underline{B}\text{ normalisateur de }\unsure{\mathscr L|\mathcal L}\arrow[rr,Leftrightarrow]
\& {} \& \text{(ii) }\mathscr L\text{ normalisateur de }\mathscr L \\
\& \text{(iii)}\arrow[u,Leftrightarrow] \& {}
\end{tikzcd}}
\diagram{\begin{tikzcd}[ampersand replacement=\&]
\operatorname{Drap}(G)\arrow[r]\arrow[d,equal]\arrow[dr]
\& \operatorname{Grass}(\unsure{\mathfrak g|\mathscr G})\arrow[d,phantom,"\cup" description]\quad\unsure{imm} \\
G/B \& G/\operatorname{Norm}_{G}(\mathscr L)
\end{tikzcd}}
\noindent\unsure{par} $g\mapsto\operatorname{int}(g).B$

% source line 6
$(\mathrm i)\Rightarrow(\mathrm{ii})$ trivial, car $\operatorname{Lie}N_{\unsure{G}}(\unsure{B|\mathscr L})=\operatorname{Norm}_{\mathfrak g}\mathscr L$

% source lines 7--10
$(\mathrm{ii})\Rightarrow(\mathrm i)$, car soit $B'$ le normalisateur, $B'\supset B$, il \unsure{a} $\unsure{=|\text{même}}$ algèbre de Lie alors $B=B'{}^0$, donc $B'$ normalise $B$ \unsure{ensemblist\textsuperscript{t}}, donc $B\unsure{\subset|=}B'$ \unsure{ensemblist\textsuperscript{t}}, donc $B=B'$

% source lines 11--19; the large outlined block is crossed out throughout.
$(\mathrm{ii})\Leftrightarrow(\mathrm{iii})$, \inserted{trivial}
\deleted{\unsure{en} \unsure{prenant} les \unsure{différentielles} \unsure{sur} les \unsure{espaces} \unsure{localisés} en $B$ resp $\mathscr L$, \unsure{on} \unsure{trouve}}
\diagram{\begin{tikzcd}[ampersand replacement=\&]
\deleted{\mathfrak g/\mathscr L}\arrow[r]\arrow[dr]
\& \deleted{\operatorname{Grass}(\mathfrak g)}\arrow[d,phantom,"\deleted{\cup}" description]\quad\deleted{\text{(}\unsure{\text{espace tangent}}\text{)}} \\
{} \& \deleted{\operatorname{Hom}(\mathscr L,\mathfrak g/\mathscr L)}
\end{tikzcd}}
\deleted{\unsure{par} $\dot x\mapsto\unsure{\mathfrak g|\varphi}D_x(\mathscr L)$, \unsure{qui} \unsure{est} \unsure{d'où} \unsure{le} \unsure{noyau} $\unsure{\to|=}\operatorname{Norm}_{\mathfrak g}(\mathscr L)/\mathscr L$. Donc $(\mathrm{ii})$ signifie que $\operatorname{Drap}(G)\to\operatorname{Grass}(\mathfrak g)$ \unsure{est} \uline{non ramifié}, \deleted{\unsure{donc}} il en résulte que \unsure{c'est} \unsure{une} \unsure{imm}, car le normalisateur de}

% source lines 20--21
Ces conditions ne sont pas tjrs vérifiées, \unsure{p.ex.} $\operatorname{SL}(2)$ en car $2$ les met en défaut.
% archive sheet number: 17
